% Point Processes Analysis Template
% Topics: Poisson processes, Hawkes processes, intensity functions, temporal clustering
% Style: Probability theory with computational simulation

\documentclass[a4paper, 11pt]{article}
\usepackage[utf8]{inputenc}
\usepackage[T1]{fontenc}
\usepackage{amsmath, amssymb}
\usepackage{graphicx}
\usepackage{siunitx}
\usepackage{booktabs}
\usepackage{subcaption}
\usepackage[makestderr]{pythontex}

% Theorem environments
\newtheorem{definition}{Definition}[section]
\newtheorem{theorem}{Theorem}[section]
\newtheorem{example}{Example}[section]
\newtheorem{remark}{Remark}[section]

\usepackage[colorlinks=false,hidelinks]{hyperref}

\title{Point Processes: Poisson and Hawkes Process Analysis}
\author{Probability and Stochastic Processes Laboratory}
\date{\today}

\begin{document}
\maketitle

\begin{abstract}
This report presents a comprehensive computational analysis of temporal point processes,
focusing on homogeneous and inhomogeneous Poisson processes and self-exciting Hawkes
processes. We implement simulation algorithms for each process type, analyze the
statistical properties of inter-arrival times and counting processes, estimate intensity
functions from observed data, and examine the clustering behavior characteristic of
Hawkes processes. The analysis demonstrates goodness-of-fit testing through residual
analysis and time-rescaling methods, with applications to modeling earthquake aftershock
sequences and financial market microstructure.
\end{abstract}

\section{Introduction}

Point processes provide a mathematical framework for modeling random events occurring
in time or space. From earthquake occurrences to financial transactions, neuron spike
trains to social media posts, point processes capture the stochastic nature of discrete
events while allowing for complex dependency structures. Unlike simple counting processes,
point processes can exhibit temporal clustering, self-excitation, and time-varying
intensity---features essential for realistic modeling of many natural and social phenomena.

The homogeneous Poisson process, with its constant event rate and memoryless property,
serves as the foundation for point process theory. However, many real-world phenomena
exhibit clustering behavior where events trigger subsequent events, violating the
independence assumption of Poisson processes. The Hawkes process, introduced by
Alan Hawkes in 1971, extends the Poisson framework by incorporating self-excitation:
each event increases the probability of future events, with this influence decaying
over time according to a kernel function.

\begin{definition}[Point Process]
A temporal point process is a random countable set of points $\{t_1, t_2, t_3, \ldots\}$
on the positive real line, where $0 < t_1 < t_2 < t_3 < \cdots$ represent event arrival
times. The counting process $N(t) = \sum_{i} \mathbf{1}_{t_i \leq t}$ gives the number
of events in $[0, t]$.
\end{definition}

In this analysis, we implement simulation algorithms for both Poisson and Hawkes
processes, examine their statistical properties through computational experiments,
and demonstrate estimation and validation techniques for fitting these models to data.

\section{Mathematical Framework}

\subsection{The Poisson Process}

The Poisson process is characterized by its intensity function $\lambda(t)$, which
specifies the instantaneous event rate at time $t$. The conditional intensity function
provides the probability of an event occurring in an infinitesimal interval.

\begin{definition}[Conditional Intensity]
The conditional intensity function $\lambda(t|\mathcal{H}_t)$ of a point process,
given the history $\mathcal{H}_t$ of all events before time $t$, satisfies:
\begin{equation}
\lambda(t|\mathcal{H}_t) = \lim_{\Delta t \to 0} \frac{\mathbb{P}(\text{event in } [t, t+\Delta t) | \mathcal{H}_t)}{\Delta t}
\end{equation}
\end{definition}

\begin{theorem}[Homogeneous Poisson Process Properties]
For a homogeneous Poisson process with constant rate $\lambda > 0$:
\begin{enumerate}
\item The number of events in interval $(s, t]$ follows $N(t) - N(s) \sim \text{Poisson}(\lambda(t-s))$
\item Inter-arrival times $\tau_i = t_i - t_{i-1}$ are i.i.d.\ $\text{Exponential}(\lambda)$
\item Given $N(T) = n$, the event times are distributed as order statistics of $n$ uniform random variables on $[0, T]$
\end{enumerate}
\end{theorem}

For inhomogeneous Poisson processes, the rate varies with time:

\begin{definition}[Inhomogeneous Poisson Process]
An inhomogeneous Poisson process with intensity function $\lambda(t)$ has:
\begin{equation}
\mathbb{P}(N(t) - N(s) = k) = \frac{\Lambda(s,t)^k}{k!} e^{-\Lambda(s,t)}
\end{equation}
where the integrated intensity is $\Lambda(s,t) = \int_s^t \lambda(u) \, du$.
\end{definition}

\subsection{The Hawkes Process}

The Hawkes process incorporates self-excitation through a memory kernel that captures
how past events influence the current intensity.

\begin{definition}[Hawkes Process]
A univariate Hawkes process has conditional intensity:
\begin{equation}
\lambda(t|\mathcal{H}_t) = \mu + \sum_{t_i < t} \phi(t - t_i)
\end{equation}
where $\mu > 0$ is the background rate and $\phi(s)$ is the excitation kernel
satisfying $\phi(s) \geq 0$ for $s > 0$ and $\phi(s) = 0$ for $s \leq 0$.
\end{definition}

\begin{theorem}[Stationarity Condition]
A Hawkes process is stationary if and only if the branching ratio satisfies:
\begin{equation}
\int_0^\infty \phi(s) \, ds < 1
\end{equation}
The branching ratio represents the expected number of offspring events triggered
by each event. Stationarity requires this to be less than 1 to prevent explosive growth.
\end{theorem}

\begin{example}[Exponential Kernel]
The most common choice is the exponential kernel $\phi(s) = \alpha e^{-\beta s}$
with $\alpha, \beta > 0$. The branching ratio is $\alpha/\beta$, and the process
is stationary when $\alpha < \beta$.
\end{example}

The mean intensity of a stationary Hawkes process is:
\begin{equation}
\mathbb{E}[\lambda(t)] = \frac{\mu}{1 - \int_0^\infty \phi(s) \, ds}
\end{equation}

\section{Simulation Algorithms}

\subsection{Poisson Process Simulation}

We implement three simulation methods for Poisson processes: the inter-arrival time
method for homogeneous processes, and thinning and inverse transform methods for
inhomogeneous processes.

\begin{pycode}
import numpy as np
import matplotlib.pyplot as plt
from scipy.optimize import minimize
from scipy.stats import kstest, expon, poisson
from scipy.special import factorial

np.random.seed(42)
plt.rcParams['text.usetex'] = True
plt.rcParams['font.family'] = 'serif'
plt.rcParams['figure.dpi'] = 150

def simulate_homogeneous_poisson(rate, T, seed=None):
    """Simulate homogeneous Poisson process via inter-arrival times."""
    if seed is not None:
        np.random.seed(seed)
    times = []
    t = 0
    while True:
        # Inter-arrival time is exponentially distributed
        tau = np.random.exponential(1.0 / rate)
        t += tau
        if t > T:
            break
        times.append(t)
    return np.array(times)

def simulate_inhomogeneous_poisson(lambda_func, lambda_max, T, seed=None):
    """Simulate inhomogeneous Poisson process via thinning algorithm."""
    if seed is not None:
        np.random.seed(seed)
    # First simulate homogeneous Poisson with rate lambda_max
    times_homo = simulate_homogeneous_poisson(lambda_max, T)
    # Then thin: keep each point with probability lambda(t)/lambda_max
    times = []
    for t in times_homo:
        if np.random.random() < lambda_func(t) / lambda_max:
            times.append(t)
    return np.array(times)

# Simulate homogeneous Poisson process
rate = 5.0  # events per unit time
T = 100.0   # observation window
events_homo = simulate_homogeneous_poisson(rate, T, seed=42)

# Compute inter-arrival times
inter_arrivals = np.diff(np.concatenate([[0], events_homo]))

# Theoretical vs empirical comparison
fig, axes = plt.subplots(2, 2, figsize=(12, 10))

# Plot 1: Event arrivals (step function)
ax = axes[0, 0]
times_plot = np.concatenate([[0], events_homo, [T]])
counts_plot = np.concatenate([[0], np.arange(1, len(events_homo)+1), [len(events_homo)]])
ax.step(times_plot, counts_plot, where='post', linewidth=1.5, color='steelblue')
ax.set_xlabel('Time $t$', fontsize=11)
ax.set_ylabel('$N(t)$', fontsize=11)
ax.set_title('Counting Process: Homogeneous Poisson ($\\lambda = 5$)', fontsize=12)
# Add theoretical expectation
t_theory = np.linspace(0, T, 100)
ax.plot(t_theory, rate * t_theory, 'r--', linewidth=1.5,
        label='$\\mathbb{E}[N(t)] = \\lambda t$', alpha=0.8)
ax.legend(fontsize=10)
ax.grid(True, alpha=0.3)

# Plot 2: Inter-arrival time distribution
ax = axes[0, 1]
ax.hist(inter_arrivals, bins=30, density=True, alpha=0.7, color='steelblue',
        edgecolor='white', label='Empirical')
x_exp = np.linspace(0, max(inter_arrivals), 100)
ax.plot(x_exp, rate * np.exp(-rate * x_exp), 'r-', linewidth=2,
        label=f'Exponential($\\lambda = {rate}$)')
ax.set_xlabel('Inter-arrival Time $\\tau$', fontsize=11)
ax.set_ylabel('Density', fontsize=11)
ax.set_title('Inter-arrival Time Distribution', fontsize=12)
ax.legend(fontsize=10)
ax.grid(True, alpha=0.3)

# Plot 3: Count distribution in fixed intervals
interval_length = 1.0
counts_per_interval = []
for i in range(int(T / interval_length)):
    start, end = i * interval_length, (i + 1) * interval_length
    count = np.sum((events_homo >= start) & (events_homo < end))
    counts_per_interval.append(count)
counts_per_interval = np.array(counts_per_interval)

ax = axes[1, 0]
unique_counts = np.arange(max(counts_per_interval) + 1)
empirical_pmf = np.array([np.sum(counts_per_interval == k) / len(counts_per_interval)
                          for k in unique_counts])
theoretical_pmf = poisson.pmf(unique_counts, rate * interval_length)
width = 0.35
ax.bar(unique_counts - width/2, empirical_pmf, width, label='Empirical',
       color='steelblue', alpha=0.7)
ax.bar(unique_counts + width/2, theoretical_pmf, width, label='Poisson(5)',
       color='coral', alpha=0.7)
ax.set_xlabel('Count $k$', fontsize=11)
ax.set_ylabel('Probability', fontsize=11)
ax.set_title('Events per Unit Interval', fontsize=12)
ax.legend(fontsize=10)
ax.grid(True, alpha=0.3, axis='y')

# Plot 4: QQ-plot for inter-arrival times
ax = axes[1, 1]
sorted_inter = np.sort(inter_arrivals)
n = len(sorted_inter)
theoretical_quantiles = expon.ppf((np.arange(1, n+1) - 0.5) / n, scale=1/rate)
ax.scatter(theoretical_quantiles, sorted_inter, alpha=0.6, s=20, color='steelblue')
max_val = max(max(theoretical_quantiles), max(sorted_inter))
ax.plot([0, max_val], [0, max_val], 'r--', linewidth=1.5, label='Perfect fit')
ax.set_xlabel('Theoretical Quantiles', fontsize=11)
ax.set_ylabel('Empirical Quantiles', fontsize=11)
ax.set_title('Q-Q Plot: Exponential Distribution', fontsize=12)
ax.legend(fontsize=10)
ax.grid(True, alpha=0.3)
ax.set_aspect('equal')

plt.tight_layout()
plt.savefig('point_processes_poisson_analysis.pdf', dpi=150, bbox_inches='tight')
plt.close()

# Perform Kolmogorov-Smirnov test
ks_stat, ks_pvalue = kstest(inter_arrivals, 'expon', args=(0, 1/rate))
\end{pycode}

\begin{figure}[htbp]
\centering
\includegraphics[width=\textwidth]{point_processes_poisson_analysis.pdf}
\caption{Comprehensive analysis of a homogeneous Poisson process with rate $\lambda = 5$
events per unit time over observation window $[0, 100]$. \textbf{Top left:} The counting
process $N(t)$ (blue steps) closely follows the theoretical expectation $\mathbb{E}[N(t)]
= \lambda t$ (red dashed), with the characteristic random fluctuations of Poisson
arrivals. \textbf{Top right:} The inter-arrival time distribution matches the
exponential density $\lambda e^{-\lambda \tau}$, confirming the memoryless property.
\textbf{Bottom left:} The number of events per unit interval follows the Poisson
distribution with parameter $\lambda = 5$. \textbf{Bottom right:} The Q-Q plot shows
excellent agreement between empirical and theoretical quantiles, validating the
exponential inter-arrival assumption.}
\end{figure}

The Kolmogorov-Smirnov test yields statistic $D = \py{f"{ks_stat:.4f}"}$ with p-value
$\py{f"{ks_pvalue:.4f}"}$, providing strong evidence that the inter-arrival times
follow the expected exponential distribution.

\subsection{Inhomogeneous Poisson Process}

Real-world event processes often exhibit time-varying intensity. We simulate an
inhomogeneous Poisson process with a sinusoidal intensity function representing
daily activity patterns.

\begin{pycode}
# Define time-varying intensity function (e.g., daily pattern)
def lambda_inhomogeneous(t):
    """Sinusoidal intensity with period 24 hours, peak at noon."""
    baseline = 3.0
    amplitude = 2.0
    period = 24.0
    return baseline + amplitude * np.sin(2 * np.pi * t / period - np.pi/2)

# Maximum intensity for thinning
lambda_max = 5.5

# Simulate for 7 days
T_days = 7 * 24  # hours
events_inhomo = simulate_inhomogeneous_poisson(lambda_inhomogeneous, lambda_max, T_days, seed=123)

fig, axes = plt.subplots(2, 2, figsize=(12, 10))

# Plot 1: Intensity function and events
ax = axes[0, 0]
t_fine = np.linspace(0, T_days, 1000)
ax.plot(t_fine, lambda_inhomogeneous(t_fine), 'r-', linewidth=2, label='$\\lambda(t)$')
ax.scatter(events_inhomo, lambda_inhomogeneous(events_inhomo), s=15, alpha=0.5,
           color='steelblue', label='Events')
ax.axhline(lambda_max, color='gray', linestyle='--', linewidth=1, label='$\\lambda_{max}$')
ax.set_xlabel('Time (hours)', fontsize=11)
ax.set_ylabel('Intensity $\\lambda(t)$', fontsize=11)
ax.set_title('Inhomogeneous Poisson Process: Sinusoidal Intensity', fontsize=12)
ax.legend(fontsize=10)
ax.grid(True, alpha=0.3)
ax.set_xlim(0, 48)  # Show first 2 days

# Plot 2: Events over time (full week)
ax = axes[0, 1]
times_plot = np.concatenate([[0], events_inhomo, [T_days]])
counts_plot = np.concatenate([[0], np.arange(1, len(events_inhomo)+1), [len(events_inhomo)]])
ax.step(times_plot, counts_plot, where='post', linewidth=1.0, color='steelblue')
# Expected count via integrated intensity
integrated_intensity = []
for t in t_fine:
    # Numerical integration of lambda(s) from 0 to t
    s_vals = np.linspace(0, t, 100)
    integral = np.trapz(lambda_inhomogeneous(s_vals), s_vals)
    integrated_intensity.append(integral)
ax.plot(t_fine, integrated_intensity, 'r--', linewidth=1.5,
        label='$\\Lambda(0,t) = \\int_0^t \\lambda(s) ds$')
ax.set_xlabel('Time (hours)', fontsize=11)
ax.set_ylabel('$N(t)$', fontsize=11)
ax.set_title('Counting Process and Integrated Intensity', fontsize=12)
ax.legend(fontsize=10)
ax.grid(True, alpha=0.3)

# Plot 3: Events by hour of day (circadian pattern)
ax = axes[1, 0]
hours_of_day = events_inhomo % 24
ax.hist(hours_of_day, bins=24, range=(0, 24), density=True, alpha=0.7,
        color='steelblue', edgecolor='white', label='Empirical')
# Theoretical density (normalized intensity)
t_hours = np.linspace(0, 24, 100)
intensity_norm = lambda_inhomogeneous(t_hours) / np.trapz(lambda_inhomogeneous(t_hours), t_hours)
ax.plot(t_hours, intensity_norm, 'r-', linewidth=2, label='Normalized $\\lambda(t)$')
ax.set_xlabel('Hour of Day', fontsize=11)
ax.set_ylabel('Density', fontsize=11)
ax.set_title('Event Distribution by Time of Day', fontsize=12)
ax.legend(fontsize=10)
ax.grid(True, alpha=0.3)
ax.set_xlim(0, 24)
ax.set_xticks([0, 6, 12, 18, 24])

# Plot 4: Time-rescaling theorem validation
ax = axes[1, 1]
# Compute rescaled inter-arrival times (should be Exp(1) if model is correct)
integrated_at_events = []
for t in events_inhomo:
    s_vals = np.linspace(0, t, 100)
    integral = np.trapz(lambda_inhomogeneous(s_vals), s_vals)
    integrated_at_events.append(integral)
integrated_at_events = np.array(integrated_at_events)
rescaled_inter = np.diff(np.concatenate([[0], integrated_at_events]))

# Compare to Exp(1)
sorted_rescaled = np.sort(rescaled_inter)
n = len(sorted_rescaled)
theoretical_quantiles = expon.ppf((np.arange(1, n+1) - 0.5) / n, scale=1)
ax.scatter(theoretical_quantiles, sorted_rescaled, alpha=0.6, s=20, color='steelblue')
max_val = max(max(theoretical_quantiles), max(sorted_rescaled))
ax.plot([0, max_val], [0, max_val], 'r--', linewidth=1.5, label='Perfect fit')
ax.set_xlabel('Theoretical Quantiles (Exp(1))', fontsize=11)
ax.set_ylabel('Rescaled Inter-arrival Times', fontsize=11)
ax.set_title('Time-Rescaling Theorem Validation', fontsize=12)
ax.legend(fontsize=10)
ax.grid(True, alpha=0.3)

plt.tight_layout()
plt.savefig('point_processes_inhomogeneous.pdf', dpi=150, bbox_inches='tight')
plt.close()

# KS test on rescaled times
ks_rescaled_stat, ks_rescaled_pvalue = kstest(rescaled_inter, 'expon', args=(0, 1))
\end{pycode}

\begin{figure}[htbp]
\centering
\includegraphics[width=\textwidth]{point_processes_inhomogeneous.pdf}
\caption{Inhomogeneous Poisson process with sinusoidal intensity $\lambda(t) = 3 + 2\sin(2\pi t/24 - \pi/2)$
modeling diurnal activity patterns. \textbf{Top left:} The intensity function oscillates
between 1 and 5 events per hour, with events (blue dots) concentrated during high-intensity
periods. \textbf{Top right:} The counting process follows the integrated intensity
$\Lambda(0,t)$, demonstrating the compensator relationship. \textbf{Bottom left:}
The distribution of events by hour of day matches the normalized intensity, with
peak activity around hour 12. \textbf{Bottom right:} The time-rescaling theorem
validation confirms model correctness: rescaled inter-arrival times follow Exp(1).}
\end{figure}

The time-rescaling theorem states that if we transform event times by the integrated
intensity $\Lambda(0, t_i)$, the transformed inter-arrival times should be i.i.d.\
$\text{Exp}(1)$ if the model is correctly specified. The Kolmogorov-Smirnov test
on rescaled times yields $D = \py{f"{ks_rescaled_stat:.4f}"}$ with p-value
$\py{f"{ks_rescaled_pvalue:.4f}"}$, validating our inhomogeneous Poisson model.

\section{Hawkes Process Simulation and Analysis}

\subsection{Simulation via Thinning}

The Hawkes process requires a modified thinning algorithm that accounts for the
self-exciting intensity. After each event, the intensity jumps and then decays
according to the excitation kernel.

\begin{pycode}
def simulate_hawkes(mu, alpha, beta, T, seed=None):
    """
    Simulate Hawkes process with exponential kernel.

    Parameters:
    - mu: background rate
    - alpha: excitation amplitude
    - beta: decay rate
    - T: observation window

    Returns array of event times.
    """
    if seed is not None:
        np.random.seed(seed)

    times = []
    t = 0

    # Intensity at any time t given history
    def intensity(t, history):
        lam = mu
        for ti in history:
            if ti < t:
                lam += alpha * np.exp(-beta * (t - ti))
        return lam

    # Upper bound on intensity (increases after each event)
    lambda_bar = mu

    while t < T:
        # Sample from exponential with current upper bound
        tau = np.random.exponential(1.0 / lambda_bar)
        t += tau

        if t > T:
            break

        # Current intensity
        lam_t = intensity(t, times)

        # Accept with probability lambda(t) / lambda_bar
        if np.random.random() < lam_t / lambda_bar:
            times.append(t)
            # Update upper bound: intensity jumps by alpha after event
            lambda_bar = lam_t + alpha
        else:
            # Update upper bound: intensity has decayed
            lambda_bar = lam_t

    return np.array(times)

# Hawkes process parameters
mu = 0.5       # background rate
alpha = 0.8    # excitation amplitude
beta = 1.0     # decay rate (branching ratio = alpha/beta = 0.8 < 1)
T_hawkes = 200

events_hawkes = simulate_hawkes(mu, alpha, beta, T_hawkes, seed=42)

# Compute intensity at each point in time
t_fine = np.linspace(0, T_hawkes, 2000)
intensity_values = np.zeros_like(t_fine)
for i, t in enumerate(t_fine):
    lam = mu
    for ti in events_hawkes:
        if ti < t:
            lam += alpha * np.exp(-beta * (t - ti))
    intensity_values[i] = lam

fig, axes = plt.subplots(2, 2, figsize=(12, 10))

# Plot 1: Intensity function with events
ax = axes[0, 0]
ax.plot(t_fine, intensity_values, 'r-', linewidth=1, alpha=0.8, label='$\\lambda(t|\\mathcal{H}_t)$')
ax.scatter(events_hawkes, np.zeros_like(events_hawkes), s=15, c='steelblue',
           marker='|', label='Events', alpha=0.7)
ax.axhline(mu, color='gray', linestyle='--', linewidth=1, label='Background $\\mu$')
ax.axhline(mu / (1 - alpha/beta), color='green', linestyle=':', linewidth=1.5,
           label=f'Mean intensity $\\mu/(1-\\alpha/\\beta) = {mu/(1-alpha/beta):.1f}$')
ax.set_xlabel('Time $t$', fontsize=11)
ax.set_ylabel('Intensity $\\lambda(t)$', fontsize=11)
ax.set_title('Self-Exciting Hawkes Process (Exponential Kernel)', fontsize=12)
ax.legend(fontsize=9, loc='upper right')
ax.grid(True, alpha=0.3)
ax.set_xlim(0, 50)  # Show first portion for clarity
ax.set_ylim(0, None)

# Plot 2: Zoomed view of clustering
ax = axes[0, 1]
# Find a cluster (period of high activity)
window_start, window_end = 20, 35
mask = (t_fine >= window_start) & (t_fine <= window_end)
events_in_window = events_hawkes[(events_hawkes >= window_start) & (events_hawkes <= window_end)]
ax.fill_between(t_fine[mask], 0, intensity_values[mask], alpha=0.3, color='red')
ax.plot(t_fine[mask], intensity_values[mask], 'r-', linewidth=2)
for ti in events_in_window:
    ax.axvline(ti, color='steelblue', alpha=0.6, linewidth=1.5)
ax.set_xlabel('Time $t$', fontsize=11)
ax.set_ylabel('Intensity $\\lambda(t)$', fontsize=11)
ax.set_title('Event Clustering: Each Event Increases Future Intensity', fontsize=12)
ax.grid(True, alpha=0.3)

# Plot 3: Counting process comparison with Poisson
ax = axes[1, 0]
# Hawkes counting process
times_plot = np.concatenate([[0], events_hawkes, [T_hawkes]])
counts_plot = np.concatenate([[0], np.arange(1, len(events_hawkes)+1), [len(events_hawkes)]])
ax.step(times_plot, counts_plot, where='post', linewidth=1.5, color='steelblue', label='Hawkes')

# Compare with Poisson with same mean rate
mean_rate = mu / (1 - alpha/beta)
events_poisson = simulate_homogeneous_poisson(mean_rate, T_hawkes, seed=123)
times_pois = np.concatenate([[0], events_poisson, [T_hawkes]])
counts_pois = np.concatenate([[0], np.arange(1, len(events_poisson)+1), [len(events_poisson)]])
ax.step(times_pois, counts_pois, where='post', linewidth=1.5, color='coral',
        alpha=0.7, label='Poisson (same mean)')

ax.plot([0, T_hawkes], [0, mean_rate * T_hawkes], 'k--', linewidth=1, label='Expected')
ax.set_xlabel('Time $t$', fontsize=11)
ax.set_ylabel('$N(t)$', fontsize=11)
ax.set_title('Counting Process: Hawkes vs Poisson', fontsize=12)
ax.legend(fontsize=10)
ax.grid(True, alpha=0.3)

# Plot 4: Inter-arrival time comparison
ax = axes[1, 1]
inter_hawkes = np.diff(np.concatenate([[0], events_hawkes]))
inter_poisson = np.diff(np.concatenate([[0], events_poisson]))

ax.hist(inter_hawkes, bins=30, density=True, alpha=0.6, color='steelblue',
        edgecolor='white', label='Hawkes')
ax.hist(inter_poisson, bins=30, density=True, alpha=0.6, color='coral',
        edgecolor='white', label='Poisson')
x_range = np.linspace(0, max(max(inter_hawkes), max(inter_poisson)), 100)
ax.plot(x_range, mean_rate * np.exp(-mean_rate * x_range), 'k--', linewidth=2,
        label='Exp($\\bar{\\lambda}$)')
ax.set_xlabel('Inter-arrival Time $\\tau$', fontsize=11)
ax.set_ylabel('Density', fontsize=11)
ax.set_title('Inter-arrival Distributions: Clustering Effect', fontsize=12)
ax.legend(fontsize=10)
ax.grid(True, alpha=0.3)

plt.tight_layout()
plt.savefig('point_processes_hawkes.pdf', dpi=150, bbox_inches='tight')
plt.close()

# Compute clustering statistics
cv_hawkes = np.std(inter_hawkes) / np.mean(inter_hawkes)  # Coefficient of variation
cv_poisson = np.std(inter_poisson) / np.mean(inter_poisson)
\end{pycode}

\begin{figure}[htbp]
\centering
\includegraphics[width=\textwidth]{point_processes_hawkes.pdf}
\caption{Self-exciting Hawkes process with parameters $\mu = 0.5$, $\alpha = 0.8$,
$\beta = 1.0$ (branching ratio $\alpha/\beta = 0.8$). \textbf{Top left:} The conditional
intensity $\lambda(t|\mathcal{H}_t)$ fluctuates around the mean intensity
$\mu/(1-\alpha/\beta) = 2.5$, with spikes following each event and exponential decay.
\textbf{Top right:} Zoomed view of a cluster showing how each event (blue lines)
increases the intensity (red curve), triggering further events in rapid succession.
\textbf{Bottom left:} Comparison of counting processes---while both reach similar
totals, the Hawkes process shows more variability due to clustering. \textbf{Bottom right:}
The Hawkes inter-arrival distribution has heavier left tail (many short intervals)
and right tail (gaps between clusters) compared to the exponential distribution.}
\end{figure}

A key diagnostic for self-excitation is the coefficient of variation (CV) of inter-arrival
times. For a Poisson process, $\text{CV} = 1$ exactly. For the Hawkes process, clustering
produces $\text{CV} > 1$ due to alternating periods of rapid events and quiet intervals.
Our simulations yield $\text{CV}_{\text{Hawkes}} = \py{f"{cv_hawkes:.3f}"}$ versus
$\text{CV}_{\text{Poisson}} = \py{f"{cv_poisson:.3f}"}$, confirming the over-dispersion
characteristic of self-exciting processes.

\subsection{Parameter Estimation}

We estimate Hawkes process parameters using maximum likelihood estimation (MLE).
The log-likelihood for a Hawkes process on $[0, T]$ is:

\begin{equation}
\log L(\mu, \alpha, \beta) = -\int_0^T \lambda(t) \, dt + \sum_{i=1}^{N(T)} \log \lambda(t_i)
\end{equation}

\begin{pycode}
def hawkes_log_likelihood(params, events, T):
    """Compute negative log-likelihood for Hawkes process."""
    mu, alpha, beta = params

    if mu <= 0 or alpha <= 0 or beta <= 0 or alpha >= beta:
        return np.inf

    n = len(events)
    if n == 0:
        return mu * T

    # Integral of intensity: int_0^T lambda(t) dt
    # = mu * T + sum_i (alpha/beta) * (1 - exp(-beta*(T - t_i)))
    integral = mu * T
    for ti in events:
        integral += (alpha / beta) * (1 - np.exp(-beta * (T - ti)))

    # Sum of log intensities
    log_intensities = 0
    for j, tj in enumerate(events):
        lam_tj = mu
        for i in range(j):
            lam_tj += alpha * np.exp(-beta * (tj - events[i]))
        if lam_tj > 0:
            log_intensities += np.log(lam_tj)
        else:
            return np.inf

    return integral - log_intensities  # Return negative for minimization

# Estimate parameters from simulated data
result = minimize(hawkes_log_likelihood, x0=[1.0, 0.5, 1.5],
                  args=(events_hawkes, T_hawkes),
                  method='L-BFGS-B',
                  bounds=[(0.01, 10), (0.01, 10), (0.01, 10)])

mu_est, alpha_est, beta_est = result.x

# Compute branching ratio
br_true = alpha / beta
br_est = alpha_est / beta_est

fig, axes = plt.subplots(1, 3, figsize=(14, 4))

# Plot 1: True vs estimated intensity
ax = axes[0]
t_fine = np.linspace(0, 50, 500)

# True intensity
intensity_true = np.zeros_like(t_fine)
for i, t in enumerate(t_fine):
    lam = mu
    for ti in events_hawkes[events_hawkes < t]:
        lam += alpha * np.exp(-beta * (t - ti))
    intensity_true[i] = lam

# Estimated intensity
intensity_est = np.zeros_like(t_fine)
for i, t in enumerate(t_fine):
    lam = mu_est
    for ti in events_hawkes[events_hawkes < t]:
        lam += alpha_est * np.exp(-beta_est * (t - ti))
    intensity_est[i] = lam

ax.plot(t_fine, intensity_true, 'b-', linewidth=2, label='True', alpha=0.7)
ax.plot(t_fine, intensity_est, 'r--', linewidth=2, label='Estimated', alpha=0.7)
ax.scatter(events_hawkes[events_hawkes < 50],
           np.zeros(np.sum(events_hawkes < 50)), s=10, c='gray', marker='|')
ax.set_xlabel('Time $t$', fontsize=11)
ax.set_ylabel('$\\lambda(t)$', fontsize=11)
ax.set_title('Intensity: True vs MLE', fontsize=12)
ax.legend(fontsize=10)
ax.grid(True, alpha=0.3)

# Plot 2: Parameter comparison
ax = axes[1]
params = ['$\\mu$', '$\\alpha$', '$\\beta$', '$\\alpha/\\beta$']
true_vals = [mu, alpha, beta, br_true]
est_vals = [mu_est, alpha_est, beta_est, br_est]

x_pos = np.arange(len(params))
width = 0.35
bars1 = ax.bar(x_pos - width/2, true_vals, width, label='True', color='steelblue', alpha=0.8)
bars2 = ax.bar(x_pos + width/2, est_vals, width, label='Estimated', color='coral', alpha=0.8)
ax.set_ylabel('Value', fontsize=11)
ax.set_title('Parameter Estimation', fontsize=12)
ax.set_xticks(x_pos)
ax.set_xticklabels(params, fontsize=11)
ax.legend(fontsize=10)
ax.grid(True, alpha=0.3, axis='y')

# Add value labels
for bar, val in zip(bars1, true_vals):
    ax.annotate(f'{val:.2f}', xy=(bar.get_x() + bar.get_width()/2, bar.get_height()),
                xytext=(0, 3), textcoords='offset points', ha='center', fontsize=9)
for bar, val in zip(bars2, est_vals):
    ax.annotate(f'{val:.2f}', xy=(bar.get_x() + bar.get_width()/2, bar.get_height()),
                xytext=(0, 3), textcoords='offset points', ha='center', fontsize=9)

# Plot 3: Residual analysis (time-rescaling)
ax = axes[2]
# Compute compensator at each event time
compensator = []
for j, tj in enumerate(events_hawkes):
    # Integral of estimated intensity from 0 to tj
    # For exponential kernel: int_0^t lambda(s) ds = mu*t + sum_{ti<t} (alpha/beta)(1 - exp(-beta(t-ti)))
    integral = mu_est * tj
    for i in range(j):
        ti = events_hawkes[i]
        integral += (alpha_est / beta_est) * (1 - np.exp(-beta_est * (tj - ti)))
    compensator.append(integral)
compensator = np.array(compensator)
rescaled_hawkes = np.diff(np.concatenate([[0], compensator]))

# Q-Q plot against Exp(1)
sorted_rescaled = np.sort(rescaled_hawkes)
n = len(sorted_rescaled)
theoretical_quantiles = expon.ppf((np.arange(1, n+1) - 0.5) / n, scale=1)
ax.scatter(theoretical_quantiles, sorted_rescaled, alpha=0.6, s=15, color='steelblue')
max_val = max(max(theoretical_quantiles), max(sorted_rescaled))
ax.plot([0, max_val], [0, max_val], 'r--', linewidth=1.5, label='Perfect fit')
ax.set_xlabel('Theoretical Quantiles', fontsize=11)
ax.set_ylabel('Rescaled Inter-arrivals', fontsize=11)
ax.set_title('Residual Q-Q Plot', fontsize=12)
ax.legend(fontsize=10)
ax.grid(True, alpha=0.3)

plt.tight_layout()
plt.savefig('point_processes_estimation.pdf', dpi=150, bbox_inches='tight')
plt.close()

# KS test on residuals
ks_hawkes_stat, ks_hawkes_pvalue = kstest(rescaled_hawkes, 'expon', args=(0, 1))
\end{pycode}

\begin{figure}[htbp]
\centering
\includegraphics[width=\textwidth]{point_processes_estimation.pdf}
\caption{Maximum likelihood estimation of Hawkes process parameters from $\py{len(events_hawkes)}$
observed events. \textbf{Left:} Comparison of true intensity (blue) and estimated
intensity (red dashed) shows excellent recovery of the temporal dynamics.
\textbf{Center:} Bar chart comparing true and estimated parameters---the MLE
successfully recovers all three parameters and the branching ratio.
\textbf{Right:} Residual Q-Q plot using time-rescaling theorem confirms good
model fit, with rescaled inter-arrivals following the Exp(1) reference distribution.}
\end{figure}

The maximum likelihood estimates are $\hat{\mu} = \py{f"{mu_est:.3f}"}$,
$\hat{\alpha} = \py{f"{alpha_est:.3f}"}$, and $\hat{\beta} = \py{f"{beta_est:.3f}"}$,
compared to true values $\mu = \py{mu}$, $\alpha = \py{alpha}$, and $\beta = \py{beta}$.
The estimated branching ratio $\hat{\alpha}/\hat{\beta} = \py{f"{br_est:.3f}"}$ closely
matches the true value $\alpha/\beta = \py{br_true}$. The Kolmogorov-Smirnov test on
the rescaled residuals yields $D = \py{f"{ks_hawkes_stat:.4f}"}$ with p-value
$\py{f"{ks_hawkes_pvalue:.4f}"}$, confirming good model fit.

\section{Applications}

\subsection{Earthquake Aftershock Modeling}

The Hawkes process is widely used in seismology to model earthquake aftershock
sequences. The Epidemic-Type Aftershock Sequence (ETAS) model uses a power-law
kernel, but the exponential Hawkes process captures the essential self-exciting
behavior.

\begin{pycode}
# Simulate earthquake sequence with realistic parameters
# Using modified parameters to represent aftershock triggering
mu_eq = 0.1      # Background seismicity rate (events/day)
alpha_eq = 1.2   # Strong aftershock triggering
beta_eq = 1.5    # Decay rate (so branching ratio = 0.8)
T_eq = 365       # One year

events_eq = simulate_hawkes(mu_eq, alpha_eq, beta_eq, T_eq, seed=456)

# Compute magnitudes (simplified: uniform distribution)
magnitudes = 2.0 + 3.0 * np.random.power(0.5, len(events_eq))  # Gutenberg-Richter-like

fig, axes = plt.subplots(2, 2, figsize=(12, 10))

# Plot 1: Event timeline with magnitudes
ax = axes[0, 0]
ax.scatter(events_eq, magnitudes, s=5*(magnitudes**2), c=magnitudes,
           cmap='YlOrRd', alpha=0.7, edgecolors='k', linewidths=0.3)
ax.set_xlabel('Time (days)', fontsize=11)
ax.set_ylabel('Magnitude', fontsize=11)
ax.set_title('Simulated Earthquake Sequence (1 Year)', fontsize=12)
ax.grid(True, alpha=0.3)
cbar = plt.colorbar(ax.collections[0], ax=ax)
cbar.set_label('Magnitude', fontsize=10)

# Plot 2: Magnitude-frequency distribution
ax = axes[0, 1]
mag_bins = np.arange(2, 5.5, 0.25)
counts, edges = np.histogram(magnitudes, bins=mag_bins)
log_counts = np.log10(counts + 0.1)
bin_centers = (edges[:-1] + edges[1:]) / 2
ax.bar(bin_centers, counts, width=0.2, alpha=0.7, color='steelblue', edgecolor='white')
ax.set_xlabel('Magnitude', fontsize=11)
ax.set_ylabel('Frequency', fontsize=11)
ax.set_title('Magnitude-Frequency Distribution', fontsize=12)
ax.set_yscale('log')
ax.grid(True, alpha=0.3, axis='y')

# Plot 3: Cumulative events showing clusters
ax = axes[1, 0]
times_plot = np.concatenate([[0], events_eq, [T_eq]])
counts_plot = np.concatenate([[0], np.arange(1, len(events_eq)+1), [len(events_eq)]])
ax.step(times_plot, counts_plot, where='post', linewidth=1.5, color='steelblue')
ax.set_xlabel('Time (days)', fontsize=11)
ax.set_ylabel('Cumulative Events', fontsize=11)
ax.set_title('Cumulative Earthquake Count', fontsize=12)
ax.grid(True, alpha=0.3)

# Mark major sequences (clusters)
# Find days with more than 5 events
day_counts = np.histogram(events_eq, bins=np.arange(0, T_eq+1, 1))[0]
high_activity = np.where(day_counts > 5)[0]
for day in high_activity[:5]:  # Mark first 5 high-activity days
    ax.axvspan(day, day+1, alpha=0.2, color='red')
ax.text(0.02, 0.98, 'Shaded: High-activity days', transform=ax.transAxes,
        fontsize=9, va='top')

# Plot 4: Inter-event times by time of year
ax = axes[1, 1]
inter_eq = np.diff(np.concatenate([[0], events_eq]))
scatter = ax.scatter(events_eq[1:], inter_eq[1:], s=20, c=magnitudes[1:],
                     cmap='YlOrRd', alpha=0.6, edgecolors='k', linewidths=0.2)
ax.set_xlabel('Time (days)', fontsize=11)
ax.set_ylabel('Inter-event Time (days)', fontsize=11)
ax.set_title('Inter-event Times Throughout the Year', fontsize=12)
ax.grid(True, alpha=0.3)
ax.set_yscale('log')
cbar = plt.colorbar(scatter, ax=ax)
cbar.set_label('Magnitude', fontsize=10)

plt.tight_layout()
plt.savefig('point_processes_earthquakes.pdf', dpi=150, bbox_inches='tight')
plt.close()

# Statistics
n_events = len(events_eq)
mean_magnitude = np.mean(magnitudes)
max_magnitude = np.max(magnitudes)
\end{pycode}

\begin{figure}[htbp]
\centering
\includegraphics[width=\textwidth]{point_processes_earthquakes.pdf}
\caption{Simulated earthquake aftershock sequence over one year using Hawkes process
with $\mu = 0.1$, $\alpha = 1.2$, $\beta = 1.5$. \textbf{Top left:} Event timeline
showing earthquake occurrences with symbol size and color proportional to magnitude.
Clustering is evident in the temporal distribution. \textbf{Top right:} The
magnitude-frequency distribution follows the expected power-law decay characteristic
of earthquake catalogs. \textbf{Bottom left:} Cumulative event count with high-activity
periods shaded in red, corresponding to aftershock sequences. \textbf{Bottom right:}
Inter-event times (log scale) showing very short times during aftershock sequences
and longer gaps during quiescent periods.}
\end{figure}

The simulated earthquake catalog contains $\py{n_events}$ events with mean magnitude
$\py{f"{mean_magnitude:.2f}"}$ and maximum magnitude $\py{f"{max_magnitude:.2f}"}$.
The clustering behavior characteristic of real seismic catalogs is clearly visible
in both the timeline and inter-event time distributions.

\section{Summary of Results}

\begin{pycode}
# Summary table
print(r'\begin{table}[htbp]')
print(r'\centering')
print(r'\caption{Summary of Point Process Analysis Results}')
print(r'\begin{tabular}{@{}lcc@{}}')
print(r'\toprule')
print(r'\textbf{Process Type} & \textbf{Parameter} & \textbf{Value} \\')
print(r'\midrule')
print(r'\multirow{3}{*}{Homogeneous Poisson} & Rate $\lambda$ & 5.0 \\')
print(f'& Events simulated & {len(events_homo)} \\\\')
print(f'& KS test p-value & {ks_pvalue:.4f} \\\\')
print(r'\midrule')
print(r'\multirow{4}{*}{Inhomogeneous Poisson} & Baseline rate & 3.0 \\')
print(r'& Amplitude & 2.0 \\')
print(f'& Events simulated & {len(events_inhomo)} \\\\')
print(f'& KS test p-value (rescaled) & {ks_rescaled_pvalue:.4f} \\\\')
print(r'\midrule')
print(r'\multirow{6}{*}{Hawkes Process} & True $\mu$ & 0.5 \\')
print(f'& Estimated $\\hat{{\\mu}}$ & {mu_est:.3f} \\\\')
print(f'& True branching ratio & {br_true:.2f} \\\\')
print(f'& Estimated branching ratio & {br_est:.3f} \\\\')
print(f'& Events simulated & {len(events_hawkes)} \\\\')
print(f'& CV of inter-arrivals & {cv_hawkes:.3f} \\\\')
print(r'\bottomrule')
print(r'\end{tabular}')
print(r'\end{table}')
\end{pycode}

\section{Conclusions}

This analysis has demonstrated the computational implementation and statistical
properties of temporal point processes, progressing from the foundational Poisson
process to the self-exciting Hawkes process. Our key findings include:

\begin{enumerate}
\item \textbf{Poisson Process Validation:} The homogeneous Poisson process with
rate $\lambda = 5$ produced $\py{len(events_homo)}$ events over $[0, 100]$, with
inter-arrival times passing the exponential distribution test (KS p-value =
$\py{f"{ks_pvalue:.4f}"}$). The event count per unit interval follows the predicted
Poisson distribution.

\item \textbf{Inhomogeneous Process Modeling:} Time-varying intensity successfully
captures periodic patterns such as circadian rhythms. The time-rescaling theorem
provides a principled validation approach, with rescaled inter-arrivals following
Exp(1) (KS p-value = $\py{f"{ks_rescaled_pvalue:.4f}"}$).

\item \textbf{Hawkes Process Self-Excitation:} The self-exciting property produces
temporal clustering with coefficient of variation $\py{f"{cv_hawkes:.3f}"}$ (compared
to 1.0 for Poisson), confirming over-dispersion due to event triggering.

\item \textbf{Parameter Estimation:} Maximum likelihood estimation successfully
recovered the true Hawkes parameters, with estimated branching ratio
$\py{f"{br_est:.3f}"}$ versus true value $\py{br_true}$.

\item \textbf{Applications:} The earthquake aftershock simulation demonstrates how
Hawkes processes capture the clustering and magnitude-frequency relationships observed
in real seismic catalogs.
\end{enumerate}

These results validate the theoretical properties of point processes and demonstrate
practical techniques for simulation, estimation, and model validation applicable to
diverse domains including seismology, finance, neuroscience, and social network analysis.

\begin{thebibliography}{20}
\bibitem{hawkes1971} A.~G.~Hawkes, ``Spectra of Some Self-Exciting and Mutually Exciting Point Processes,'' \textit{Biometrika}, vol.~58, no.~1, pp.~83--90, 1971.

\bibitem{hawkes1974} A.~G.~Hawkes and D.~Oakes, ``A Cluster Process Representation of a Self-Exciting Process,'' \textit{Journal of Applied Probability}, vol.~11, no.~3, pp.~493--503, 1974.

\bibitem{ogata1988} Y.~Ogata, ``Statistical Models for Earthquake Occurrences and Residual Analysis for Point Processes,'' \textit{Journal of the American Statistical Association}, vol.~83, no.~401, pp.~9--27, 1988.

\bibitem{ogata1998} Y.~Ogata, ``Space-Time Point-Process Models for Earthquake Occurrences,'' \textit{Annals of the Institute of Statistical Mathematics}, vol.~50, no.~2, pp.~379--402, 1998.

\bibitem{daley2003} D.~J.~Daley and D.~Vere-Jones, \textit{An Introduction to the Theory of Point Processes, Volume I: Elementary Theory and Methods}, 2nd ed. New York: Springer, 2003.

\bibitem{daley2008} D.~J.~Daley and D.~Vere-Jones, \textit{An Introduction to the Theory of Point Processes, Volume II: General Theory and Structure}, 2nd ed. New York: Springer, 2008.

\bibitem{embrechts2011} P.~Embrechts, T.~Liniger, and L.~Lin, ``Multivariate Hawkes Processes: An Application to Financial Data,'' \textit{Journal of Applied Probability}, vol.~48A, pp.~367--378, 2011.

\bibitem{bacry2015} E.~Bacry, I.~Mastromatteo, and J.-F.~Muzy, ``Hawkes Processes in Finance,'' \textit{Market Microstructure and Liquidity}, vol.~1, no.~1, p.~1550005, 2015.

\bibitem{lewis1979} P.~A.~W.~Lewis and G.~S.~Shedler, ``Simulation of Nonhomogeneous Poisson Processes by Thinning,'' \textit{Naval Research Logistics Quarterly}, vol.~26, no.~3, pp.~403--413, 1979.

\bibitem{ozaki1979} T.~Ozaki, ``Maximum Likelihood Estimation of Hawkes' Self-Exciting Point Processes,'' \textit{Annals of the Institute of Statistical Mathematics}, vol.~31, no.~1, pp.~145--155, 1979.

\bibitem{brown2002} E.~N.~Brown, R.~Barbieri, V.~Ventura, R.~E.~Kass, and L.~M.~Frank, ``The Time-Rescaling Theorem and Its Application to Neural Spike Train Data Analysis,'' \textit{Neural Computation}, vol.~14, no.~2, pp.~325--346, 2002.

\bibitem{rubin1972} I.~Rubin, ``Regular Point Processes and Their Detection,'' \textit{IEEE Transactions on Information Theory}, vol.~18, no.~5, pp.~547--557, 1972.

\bibitem{kingman1992} J.~F.~C.~Kingman, \textit{Poisson Processes}. Oxford: Clarendon Press, 1992.

\bibitem{moller2003} J.~M{\o}ller and R.~P.~Waagepetersen, \textit{Statistical Inference and Simulation for Spatial Point Processes}. Boca Raton: Chapman \& Hall/CRC, 2003.

\bibitem{reinhart2018} A.~Reinhart, ``A Review of Self-Exciting Spatio-Temporal Point Processes and Their Applications,'' \textit{Statistical Science}, vol.~33, no.~3, pp.~299--318, 2018.

\bibitem{rasmussen2011} J.~G.~Rasmussen, ``Temporal Point Processes: The Conditional Intensity Function,'' \textit{arXiv preprint arXiv:1806.00221}, 2018.

\bibitem{zhuang2004} J.~Zhuang, Y.~Ogata, and D.~Vere-Jones, ``Stochastic Declustering of Space-Time Earthquake Occurrences,'' \textit{Journal of the American Statistical Association}, vol.~99, no.~465, pp.~369--379, 2004.

\bibitem{helmstetter2002} A.~Helmstetter and D.~Sornette, ``Subcritical and Supercritical Regimes in Epidemic Models of Earthquake Aftershocks,'' \textit{Journal of Geophysical Research}, vol.~107, no.~B10, p.~2237, 2002.

\bibitem{laub2015} P.~J.~Laub, T.~Taimre, and P.~K.~Pollett, ``Hawkes Processes,'' \textit{arXiv preprint arXiv:1507.02822}, 2015.

\bibitem{crane2008} R.~Crane and D.~Sornette, ``Robust Dynamic Classes Revealed by Measuring the Response Function of a Social System,'' \textit{Proceedings of the National Academy of Sciences}, vol.~105, no.~41, pp.~15649--15653, 2008.
\end{thebibliography}

\end{document}
